\chapter{Rust for Low Latency}\label{ch:ll:rust-for-low-latency}

The Rust twin of this book's feed decoder decodes a message in about eleven nanoseconds on this laptop, within a fifth
of the C++ original from one run to the next. Both check every length before reading, neither allocates per message, and the
Rust version was written without a single \texttt{unsafe} block. What Rust changes is not speed but where mistakes are
caught: a \omterm{def:ll:rust-for-low-latency:race}{data race} or a dangling reference is a compile error, an out-of-bounds index a controlled panic instead of a
silent read of someone else's memory. This chapter looks at Rust on the hot path the way chapters~6 to~8 looked at
C++: which abstractions really cost nothing, where the checks are and when the compiler removes them, how to write
\texttt{unsafe} code that stays \omterm{def:ll:rust-for-low-latency:sound}{sound}, how to control allocation and threads, and where Rust's async machinery belongs
in a trading system and where it does not.

\section{Zero-cost abstractions, checked}

\begin{definition}[Zero-cost abstraction]\label{def:ll:rust-for-low-latency:zero}
A \emph{zero-cost abstraction}\index{zero-cost abstraction} compiles to the code a programmer would have written by hand
without it: the abstraction's convenience is resolved at compile time and leaves no run-time trace. The formulation comes
from C++: ``what you don't use, you don't pay for; and what you do use, you couldn't hand code any better''
(Stroustrup).
\end{definition}

Iterators, closures, generics and \texttt{Option} are designed to be zero-cost in Rust: they are monomorphised like C++
templates and inlined. The claim is checked the same way as in chapter~7, in the assembly. Safe Rust adds one thing C++
does not: every slice index is checked against the slice's length.

\begin{definition}[Bounds-check elimination]\label{def:ll:rust-for-low-latency:bce}
\emph{Bounds-check elimination}\index{bounds-check elimination} is the compiler's removal of an index check it can prove
always passes: a loop counter bounded by the slice's own length, an index already compared, an iterator that yields only
valid positions.
\end{definition}

\Cref{lst:ll:rust-for-low-latency:counted} is the assembly of a loop \texttt{for i in 0..v.len() \{ s += v[i] \}} compiled
by \texttt{rustc -O}: no comparison against the length and no call to the panic handler remain, and LLVM has vectorised the
loop four elements at a time. \Cref{lst:ll:rust-for-low-latency:checked} reads one element at an index the compiler knows
nothing about: a comparison, a predicted branch, and the panic path out of line. When the check cannot be removed it costs
little, as long as it is predicted and the load it guards is the slow part: in
\cref{fig:ll:rust-for-low-latency:twins} a gather of 65\,536 elements through a table of scattered indices costs the same,
within noise, with and without the check, in both languages.

\omcode{code/low-latency/09-rust-for-low-latency/rust/asm/sum_counted.s}{1}{40}{A counted loop over a slice, rustc 1.97 -O: no bounds check, vectorised.}\label{lst:ll:rust-for-low-latency:counted}

\omcode{code/low-latency/09-rust-for-low-latency/rust/asm/get_checked.s}{1}{12}{One indexed read at an unknown index: a comparison and a branch to the panic handler.}\label{lst:ll:rust-for-low-latency:checked}

\begin{omfigure}
\begin{tikzpicture}
\begin{axis}[name=dec,width=4.6cm,height=5.4cm,ybar,bar width=12pt,ymin=0,symbolic x coords={decode},xtick=data,
  enlarge x limits=0.8,title={\small ns per message},tick label style={font=\footnotesize},label style={font=\small},
  grid=major,grid style={gray!25},table/col sep=comma]
\addplot[fill=omDef!55,draw=omDef] table[x=task,y=cpp]{figdata/low-latency/09-rust-for-low-latency/measured_decode.csv};
\addplot[fill=omThm!45,draw=omThm] table[x=task,y=rust]{figdata/low-latency/09-rust-for-low-latency/measured_decode.csv};
\end{axis}
\begin{axis}[at={(dec.east)},anchor=west,xshift=1.5cm,width=8.6cm,height=5.4cm,ybar,bar width=10pt,ymin=0,
  symbolic x coords={counted loop,gather checked,gather unchecked},xtick=data,enlarge x limits=0.25,
  xticklabels={{counted\\loop},{gather\\checked},{gather\\unchecked}},xticklabel style={align=center},
  title={\small ns per element},tick label style={font=\footnotesize},label style={font=\small},grid=major,grid style={gray!25},
  legend columns=2,legend style={font=\footnotesize,at={(0.3,-0.33)},anchor=north,draw=none,fill=none,
  /tikz/every even column/.append style={column sep=8pt}},table/col sep=comma]
\addplot[fill=omDef!55,draw=omDef] table[x=task,y=cpp]{figdata/low-latency/09-rust-for-low-latency/measured_loops.csv};
\addplot[fill=omThm!45,draw=omThm] table[x=task,y=rust]{figdata/low-latency/09-rust-for-low-latency/measured_loops.csv};
\legend{C++ (g++ -O2),Rust (rustc -O)}
\end{axis}
\end{tikzpicture}

\omcaption{The C++ and Rust twins: decoding Book~1's sample with each language's feed decoder (left), and summing 65\,536
integers in order, and through a table of scattered indices with and without bounds checks (right). Measured on a laptop
(Intel Core Ultra 7 155H) under WSL2, no isolated cores, pinned to one CPU. Data: \texttt{bench\_rust.py}.}\label{fig:ll:rust-for-low-latency:twins}
\end{omfigure}

\section{Unsafe code and its boundaries}

Rust's guarantees hold for safe code; \texttt{unsafe} blocks are where the programmer takes over the compiler's proof
obligations: dereferencing raw pointers, calling C functions, skipping bounds checks, implementing lock-free structures
(chapter~11).

\begin{definition}[Sound abstraction]\label{def:ll:rust-for-low-latency:sound}
An abstraction built on \texttt{unsafe} code is \emph{sound}\index{sound abstraction} if no sequence of calls to its safe
interface, however unusual, can cause \omterm{def:ll:cpp-for-latency-iii-the-compiler:ub}{undefined behaviour}: the invariants the unsafe code relies on are established and
maintained by the abstraction itself, not assumed of its callers.
\end{definition}

The discipline is to keep the unsafe region small and to write down, next to each block, why it is \omterm{def:ll:rust-for-low-latency:sound}{sound}. The gather of
\cref{lst:ll:rust-for-low-latency:sound} skips the per-element check only after the indices have been validated once, and
refuses a slice whose length is not the one they were validated against: however it is called, it cannot read out of
bounds. A function whose safety depends on its caller is itself marked \texttt{unsafe}, with a \texttt{\# Safety} section
stating the contract; \texttt{clippy} rejects a public function that dereferences a raw pointer argument without it. The
same questions apply to C++ code, where every line is unsafe; Rust makes them local.

\omcode{code/low-latency/09-rust-for-low-latency/rust/src/lib.rs}{58}{72}{A sound interface over an unchecked gather.}\label{lst:ll:rust-for-low-latency:sound}

\section{Controlling allocation}

\begin{definition}[Global allocator]\label{def:ll:rust-for-low-latency:alloc}
The \emph{global allocator}\index{global allocator} of a Rust program is the allocator behind \texttt{Box}, \texttt{Vec},
\texttt{String} and every standard collection; a program can replace it with any type implementing \texttt{GlobalAlloc},
declared with \texttt{\#[global\_allocator]}, for example one that counts allocations in tests.
\end{definition}

Rust's standard types allocate in slightly different places from C++'s, and the differences matter on a hot path. A
\texttt{String} has no inline buffer, so even a four-letter symbol allocates; a \texttt{Vec} reserves room for four small
elements at its first push, so three fills cost one allocation instead of three; \texttt{HashMap::entry} moves the key into
the map instead of copying it. Counted with \texttt{firm\_arena::CountingAlloc}, the Rust version of chapter~6's first-draft
handler (\cref{lst:ll:rust-for-low-latency:naive}) allocates five times per message where the C++ one allocated nine, and the
fixed version none. The remedy is the same as in C++: fixed arrays, pools of indices, byte arrays for identifiers, no boxed
closures on the hot path, and a test that asserts the count.

\omcode{code/low-latency/09-rust-for-low-latency/rust/src/lib.rs}{89}{102}{Chapter 6's first draft in Rust: five allocations per message.}\label{lst:ll:rust-for-low-latency:naive}

\section{Threads, pinning and what the type system guarantees}

\begin{definition}[Data race]\label{def:ll:rust-for-low-latency:race}
A \emph{data race}\index{data race} occurs when two threads access the same memory location concurrently, at least one
access is a write, and at least one is not atomic. In C++ it is \omterm{def:ll:cpp-for-latency-iii-the-compiler:ub}{undefined behaviour}; in safe Rust it cannot be written:
the traits \texttt{Send} and \texttt{Sync}, which mark the types that may be moved or shared between threads, and the rule
that a mutable reference is exclusive, are checked when compiling.
\end{definition}

The chapter's crate carries a documentation test marked \texttt{compile\_fail}: a thread that increments an element of a
vector while its owner also writes it. The test passes only if the compiler rejects the program, which it does, because
the spawned thread borrows the vector mutably while it is still in use. What the type system does not provide is
placement: the standard library cannot pin a thread to a CPU. \texttt{firm.affinity} declares the two C functions it needs,
calls them inside \texttt{unsafe} blocks with the reason stated, and verifies the result with \texttt{sched\_getcpu}
(\cref{lst:ll:rust-for-low-latency:pin}); its C++ twin does the same with \texttt{pthread\_setaffinity\_np}. Both read the
role-to-CPU plan that \texttt{firm.coreplan} writes (chapter~4).

\omcode{code/firm/affinity/rust/src/lib.rs}{36}{45}{Pinning the calling thread from Rust through the C library.}\label{lst:ll:rust-for-low-latency:pin}

\section{Where async does and does not belong}

\begin{definition}[Async runtime]\label{def:ll:rust-for-low-latency:async}
An \emph{async runtime}\index{async runtime} runs many tasks written as \texttt{async} functions on a few threads: a task
that would wait (for a socket, a timer) returns control to the runtime's scheduler, which resumes it when the event
arrives, possibly on another thread.
\end{definition}

Async code is how Rust serves many slow connections cheaply, and a trading firm has many: venue REST interfaces, websocket
market data from dozens of crypto venues (chapter~17), drop-copy sessions, internal services. The hot path is the opposite
case: one busy thread per role, pinned, spinning on its input, never yielding. A scheduler that may move a task between
threads, or delay it behind another, adds exactly the \omterm{def:ll:where-latency-comes-from:tail}{jitter} the hot path exists to avoid. The authors of the most
widely used runtime say so themselves: it is designed for applications whose tasks spend most of their time waiting for
input and output, not for speeding up computation. The usual split is a synchronous, pinned hot path, fed and drained by
lock-free rings (chapter~12), with async services around it for everything that can wait.

\section{Tutorial: the Rust twin, measured}\label{tut:ll:rust-for-low-latency:tut}

\begin{tutorial}
\textbf{Goal.} Check Rust's abstractions in the assembly, count its allocations, pin a thread, and compare the twins.
\textbf{End state:} the two listings of assembly, the allocation counts, and \cref{fig:ll:rust-for-low-latency:twins}.
\begin{tutsteps}
\item \textbf{Assembly.} \texttt{python ll\_rustasm.py} compiles the crate with \texttt{rustc -O --emit asm} and extracts
\texttt{sum\_counted} and \texttt{get\_checked}, exported with \texttt{\#[no\_mangle]} so that their names are stable.
\item \textbf{Soundness and a compile-time race.} \texttt{cargo test} runs the unit tests and the \texttt{compile\_fail}
documentation test; clippy runs with \texttt{-D warnings}.
\item \textbf{Allocations.} The test installs \texttt{CountingAlloc} as the \omterm{def:ll:rust-for-low-latency:alloc}{global allocator} and asserts five allocations
per message for the draft and none for the fixed handler.
\item \textbf{Pinning.} \texttt{firm\_affinity::apply(plan, role)} pins a thread and verifies its CPU.
\item \textbf{Measure} with \texttt{python bench\_rust.py}: it builds the crate's benchmark with \texttt{cargo build
--release} and the C++ twin with g++ and runs both pinned.
\end{tutsteps}
\textbf{What to change next.} Replace the fixed handler's linear search with a small open-addressing table and keep the
allocation count at zero; compile the C++ twin with \texttt{-O3} and see which language's counted loop is faster then.
\end{tutorial}

\section{Build: thread placement by role}\label{bld:ll:rust-for-low-latency:affinity}

\begin{build}
\bfield{Purpose} Every hot thread of Part~IV is pinned to the CPU the placement plan gives its role, and says so; the
measurements of chapter~26 record where each thread ran.
\bfield{Interface} Rust \texttt{firm\_affinity}: \texttt{read\_plan(text)}, \texttt{read\_plan\_file(path)},
\texttt{pin\_current(cpu)}, \texttt{current\_cpu()}, \texttt{allowed()}, \texttt{apply(plan, role)}. C++20
\texttt{firm::affinity}: \texttt{read\_plan(path)}, \texttt{pin\_current}, \texttt{current\_cpu}, \texttt{name\_current},
\texttt{apply}. Python \texttt{firm\_coreplan.write\_plan} and \texttt{read\_plan} produce and read the plan file.
\bfield{Rules} A thread is pinned before it allocates or touches its data; \texttt{apply} fails, rather than continuing
unpinned, if the role is missing, the kernel refuses, or the thread is not on its CPU afterwards; every \texttt{unsafe}
call carries its safety argument.
\bfield{Acceptance tests} \texttt{code/firm/affinity/}: the example plan read in both languages; a thread pinned to an
allowed CPU and verified; a missing role refused. \texttt{code/firm/coreplan/}: the plan file round trip.
\bfield{Stretch} Set the thread's name for \texttt{top} and \texttt{perf}; refuse to start when the plan puts two hot roles
on siblings.
\end{build}

\begin{omsources}
\item B.~Stroustrup, ``Abstraction and the C++ machine model'', 2005.
\item The Rustonomicon, ``Send and Sync''; The Rust Reference, ``Unsafety''.
\item Tokio, tutorial, ``When not to use Tokio''.
\end{omsources}

\section{Exercises}

\begin{exercise}[$\star$]\label{exo:ll:rust-for-low-latency:1}
From \cref{fig:ll:rust-for-low-latency:twins}, by what factor is the Rust decoder slower or faster than the C++ one, and
what does that cost at 5 million messages a second?
\end{exercise}

\begin{exercise}[$\star$]\label{exo:ll:rust-for-low-latency:2}
Which of these allocate in Rust: \texttt{\textquotedbl{}ESZ6\textquotedbl{}.to\_string()}; \texttt{Vec::<u64>::new()}; the first
\texttt{push} into it; \texttt{[0u8; 24]}; \texttt{Box::new(3)}?
\end{exercise}

\begin{exercise}[$\star$]\label{exo:ll:rust-for-low-latency:3}
Why does the Rust draft handler allocate five times where the C++ one allocated nine? Account for the difference item by
item.
\end{exercise}

\begin{exercise}[$\star\star$]\label{exo:ll:rust-for-low-latency:4}
Read \cref{lst:ll:rust-for-low-latency:counted}: how many elements does the vector loop add per iteration, and what handles
lengths that are not a multiple of four?
\end{exercise}

\begin{exercise}[$\star\star$]\label{exo:ll:rust-for-low-latency:5}
Why does removing the bounds check from the gather not make it faster here? When would it?
\end{exercise}

\begin{exercise}[$\star\star$]\label{exo:ll:rust-for-low-latency:6}
Is \texttt{CheckedIndex} still \omterm{def:ll:rust-for-low-latency:sound}{sound} if \texttt{sum} does not compare \texttt{v.len()} with \texttt{len}? Give a calling
sequence that would break it.
\end{exercise}

\begin{exercise}[$\star\star\star$]\label{exo:ll:rust-for-low-latency:7}
\emph{Coding.} Write the Rust decoder's hot loop over \texttt{decode(\&data).flatten()} with an explicit \texttt{match} on
the result instead, count allocations with \texttt{CountingAlloc}, and check the assembly for panics. What differs?
\end{exercise}

\begin{exercise}[$\star\star\star$]\label{exo:ll:rust-for-low-latency:8}
\emph{Find the flaw.} ``Our market-data handler is an async task on a multi-threaded runtime so that it scales with the
number of cores.''
\end{exercise}

\section{Problem: The Port}

\begin{problem}[{Weekend problem --- rewriting the feed path in Rust}]\label{pb:ll:rust-for-low-latency:1}
A firm considers porting its feed handler and order-book builder from C++ to Rust. The team measures the twins
(\cref{fig:ll:rust-for-low-latency:twins}), counts allocations, and reads the assembly.

\medskip\noindent\textbf{Part I --- Speed.}
\begin{enumerate}
\item What does each decoder cost per message, and what is the ratio?
\item Why is the Rust counted loop faster than the C++ one with the flags used?
\item What would make the C++ loop match it?
\item What do the gather measurements say about bounds checks on this workload?
\end{enumerate}

\medskip\noindent\textbf{Part II --- Allocation.}
\begin{enumerate}[resume]
\item How many allocations per message does each draft make, and each fixed handler?
\item Which Rust type allocates where the C++ one did not?
\item How does the test enforce zero allocations in Rust?
\item What happens if a dependency allocates on the hot path without the team knowing?
\end{enumerate}

\medskip\noindent\textbf{Part III --- Safety.}
\begin{enumerate}[resume]
\item Which class of bugs does the port remove at compile time?
\item Where does \texttt{unsafe} remain in the firm's Rust code, and why?
\item What does a \omterm{def:ll:rust-for-low-latency:sound}{sound abstraction} promise its callers?
\item What does a panic on the hot path mean in production, and how should it be handled?
\end{enumerate}

\medskip\noindent\textbf{Part IV --- The decision.}
\begin{enumerate}[resume]
\item State the \emph{named result}: the Rust-to-C++ ratio of time per message for the decoder, and the allocations per
message of each draft and fixed handler.
\item Is the speed difference a reason not to port?
\item What would you port first, and why?
\item Where does async Rust belong in the firm's systems?
\item How do the two languages share the firm's placement plan and fixtures?
\item What would the team lose by porting, apart from speed?
\item How would you keep both implementations equivalent during the transition?
\item In one sentence: what does Rust change for a low-latency team?
\end{enumerate}
\end{problem}

\section{Interview questions}

\begin{interviewq}[$\star$ \iqroles{developer}]\label{iq:ll:rust-for-low-latency:1}
What does ``\omterm{def:ll:rust-for-low-latency:zero}{zero-cost abstraction}'' mean, and how would you verify the claim for an iterator chain?
\end{interviewq}

\begin{interviewq}[$\star\star$ \iqroles{developer}]\label{iq:ll:rust-for-low-latency:2}
Do bounds checks make Rust slow? When do they cost, and how do you remove them safely?
\end{interviewq}

\begin{interviewq}[$\star\star$ \iqroles{developer}]\label{iq:ll:rust-for-low-latency:3}
What makes an \texttt{unsafe} abstraction \omterm{def:ll:rust-for-low-latency:sound}{sound}?
\end{interviewq}

\begin{interviewq}[$\star\star$ \iqroles{developer}]\label{iq:ll:rust-for-low-latency:4}
How does Rust prevent \omterm{def:ll:rust-for-low-latency:race}{data races}, and what does it not prevent?
\end{interviewq}

\begin{interviewq}[$\star\star$ \iqroles{developer}]\label{iq:ll:rust-for-low-latency:5}
Would you use async Rust on a trading hot path? Why?
\end{interviewq}

\begin{interviewq}[$\star\star\star$ \iqroles{developer}]\label{iq:ll:rust-for-low-latency:6}
How would you prove that a Rust component makes no heap allocation after start-up, including inside its dependencies?
\end{interviewq}
